\documentclass[12pt]{article}
% Préambule commun à tous les documents extraits de « La Revue CAPES »
\usepackage[T1]{fontenc}
\usepackage[utf8]{inputenc}
\usepackage{lmodern}
\usepackage{amsmath,amssymb,amsthm,amsfonts,mathrsfs}
\usepackage{setspace}
\usepackage{listings}
\usepackage{pgfplots}
\pgfplotsset{compat=1.18}
\usepackage{longtable}
\usepackage{adjustbox}
\usepackage{lettrine}
\usepackage{tkz-tab}
\usepackage{changepage}
\usepackage{pdflscape}
\usepackage{graphicx}
\usepackage{tikz}
\usepackage{xcolor}
\usepackage[a4paper,top=2cm,bottom=2.5cm,left=2.5cm,right=2cm]{geometry}
\usepackage{fancyhdr}
\usepackage{draftwatermark}
\SetWatermarkText{La RevueCapes}
\SetWatermarkLightness{0.9}
\SetWatermarkScale{2}
\usepackage{hyperref}
\hypersetup{colorlinks=true,linkcolor=blue!50!black,urlcolor=blue!50!black}

\definecolor{revue}{RGB}{20,60,120}

\pagestyle{fancy}
\fancyhf{}
\renewcommand{\headrulewidth}{0.4pt}
\setlength{\headheight}{15pt}

% \entete{Matière}{Titre}{Type de document}
\newcommand{\entete}[3]{%
  \fancyhead[L]{\small\textsc{La Revue CAPES} -- #1}%
  \fancyhead[R]{\small #2}%
  \fancyfoot[C]{\thepage}%
  \thispagestyle{plain}%
  \begin{center}
    {\color{revue}\rule{\textwidth}{1.2pt}}\\[4pt]
    {\small\textsc{Concours directs CAP-CEG / CAPES -- Burkina Faso}}\\[6pt]
    {\LARGE\bfseries\color{revue} #2}\\[6pt]
    {\large\itshape #1 -- #3}\\[2pt]
    {\color{revue}\rule{\textwidth}{1.2pt}}
  \end{center}
  \vspace{0.5cm}%
}

% Encadré pour les remarques ajoutées dans les corrigés
\newcommand{\remarque}[1]{\par\medskip\noindent\fcolorbox{revue}{revue!6}{\parbox{\dimexpr\linewidth-2\fboxsep-2\fboxrule}{\textbf{\color{revue}Remarque.} #1}}\par\medskip}


\begin{document}
\entete{Mathématiques}{Ancien sujet CAPES -- Session 2018}{Énoncé}
\begin{spacing}{1.5}

\textbf{Exercice 1 : (5pts)}

On donne $A=\begin{pmatrix}
2 & 2& 2 \\ 
-1 & -1 & 1 \\ 
3 & 3 & 1
\end{pmatrix} $

1.a. Calculer les valeurs propres de $A$ et en déduire que A est diagonalisable.

1.b. Déterminer une matrice inversible $P$ et une matrice diagonale $D$ tel que $A=PDP^{-1}$

2.a. Montrer que 
\[
\forall {n\in\mathbb{N}^*}, A^n = PD^nP^{-1}
\]
et en déduire $A^n$ en fonction de $n$.

2.b. Appliquer ce résultat au  calcul de $A^4$.

\textbf{Exercice 2 : (5pts)}

Soit l'application : 
\begin{align*}
f:\mathbb{R}^3 &\longrightarrow \mathbb{R}^3\\
(x,y,z) &\longmapsto f(x,y,z)=(-x+y+z, -6x+4y+2z, 3x-y+z)
\end{align*} 

1. Écrire la matrice de $f$ dans la base canonique de $\mathbb{R}^3$.

2. Montrer l'égalité $f\circ f= 2f$. En déduire que si $v\in Imf$, alors $f(v)=2v$.

3. Montrer que les sous-espaces vectoriels $Kerf$ et $Imf$ sont des sous-espaces supplémentaires de $\mathbb{R}^3$

4. Trouver une base $(e_1,e_2,e_3)$ de $\mathbb{R}^3$ dont le premier vecteur appartient à $Kerf$ et les derniers à $Imf$. Quelle est la matrice de $f$ dans la base $(e_1,e_2,e_3)$ ?

5. Trouver une équation de $Imf$.

\textbf{Exercice 3 : (5pts)}

On pose pour $(x,y)\in\Omega=\mathbb{R}\times]0,+\infty[$,  $f(x,y)=y(x^2 + ln^2{(y)})$.

1.a. Calculer $\frac{\partial{f}}{\partial{x}}(x,y)$ et trouver les $(x,y)\in\Omega$ tels que $\frac{\partial{f}}{\partial{x}}(x,y) = 0$.

1.b. Calculer $\frac{\partial{f}}{\partial{y}}(x,y)$ et en déduire les points critiques de $f$.

2.a. En déduire les extrema relatifs de $f$.

2.b. Montrer que si $(x,y)\in\Omega$ et $(x,y)\neq(0,1)$, on a $f(x,y)>0$. Qu'en déduire pour $(0,1)$?

\textbf{Exercice 4 : (5pts)}

a. Soit $\alpha\in \mathbb{R}$. Résoudre l'équation $(E_1):y^{\prime\prime} -2y^{\prime} +(1-\alpha^2)y=e^x$.

b. En déduire la solution de l'équation différentielle $(E_2):y^{\prime\prime} -2y^{\prime}+y=e^x$ satisfaisant au conditions initiales $y(0)=1$ et $y^\prime(0)=3$.

\medskip

\end{spacing}
\end{document}
